\documentclass[a4paper]{arrowhead}

\usepackage[yyyymmdd]{datetime}
\usepackage{etoolbox}
\usepackage[utf8]{inputenc}
\usepackage[T1]{fontenc}
\usepackage{multirow}
\usepackage{longtable}
\usepackage{tikz}
\usetikzlibrary{positioning,shapes.multipart,arrows.meta}

\renewcommand{\dateseparator}{-}
\setlength{\emergencystretch}{3em}
\lstset{language=HTTP,breakatwhitespace=false,columns=fullflexible}

\newcommand{\fparam}[1]{\textit{\textcolor{ArrowheadBlue}{#1}}}

%% Special references
\newcommand{\mref}[2][]{{\textcolor{ArrowheadPurple}{\hyperref[sec:model:#2]{\ifstrempty{#1}{#2}{#1}}}}}
\newcommand{\oref}[1]{{\textcolor{ArrowheadBlue}{\hyperref[sec:operations:#1]{#1}}}}
\newcommand{\pdef}[1]{{\textcolor{ArrowheadGrey}{#1 \label{sec:model:primitives:#1} \label{sec:model:primitives:#1s}}}}
\newcommand{\pref}[1]{{\textcolor{ArrowheadGrey}{\hyperref[sec:model:primitives:#1]{#1}}}}

% as template_idd.tex; the input and output arguments are passed already formatted (\mref or \pref)
\newrobustcmd\osubsection[5]{
  \addtocounter{subsection}{1}
  \addcontentsline{toc}{subsection}{\protect\numberline{\thesubsection}operation \textcolor{ArrowheadBlue}{#1}}
  \subsection*{
    \thesubsection\quad
    #2 \textcolor{ArrowheadPurple}{#3} \\
    \small
    \hspace*{0.075\textwidth}\begin{minipage}{0.1\textwidth}
      \vspace*{1mm}
      Operation: \\
      \notblank{#4}{Input: \\}{}
      \notblank{#5}{Output: \\}{}
    \end{minipage}
    \begin{minipage}{0.825\textwidth}
      \vspace*{1mm}
      \textcolor{ArrowheadBlue}{#1} \\
      \notblank{#4}{#4 \\}{}
      \notblank{#5}{#5 \\}{}
    \end{minipage}
  }
  \label{sec:operations:#1}
}
\newrobustcmd\msubsection[2]{
  \addtocounter{subsection}{1}
  \addcontentsline{toc}{subsection}{\protect\numberline{\thesubsection}#1 \textcolor{ArrowheadPurple}{#2}}
  \subsection*{\thesubsection\quad#1 \textcolor{ArrowheadPurple}{#2}}
  \label{sec:model:#2} \label{sec:model:#2s}
}
%%

\usepackage{xurl}
\begin{document}
\sloppy

%% Arrowhead Document Properties
\ArrowheadTitle{globalServiceDiscovery HTTP/TLS/JSON}
\ArrowheadServiceID{globalServiceDiscovery}
\ArrowheadType{Interface Design Description}
\ArrowheadTypeShort{IDD}
\ArrowheadVersion{5.3.0}
\ArrowheadDate{2026-11-09}
\ArrowheadAuthor{Jerker Delsing}
\ArrowheadStatus{REVIEW DRAFT}
\ArrowheadContact{jerker.delsing@ltu.se}
\ArrowheadFooter{\href{www.arrowhead.eu}{www.arrowhead.eu}}
\ArrowheadSetup
%%

%% Front Page
\begin{center}
  \vspace*{1cm}
  \huge{\arrowtitle}

  \vspace*{0.2cm}
  \LARGE{\arrowtype}
  \vspace*{1cm}

  \Large{Service ID: \textit{"\arrowid"}}

  \vspace*{\fill}

  % Front Page Image

  \vspace*{1cm}
  \vspace*{\fill}

  % Front Page Abstract
  \begin{abstract}
    This document describes the HTTP/1.1 protocol with TLS payload security and JSON payload encoding variant of the \texttt{global\allowbreak{}Service\allowbreak{}Discovery} service of the Gatekeeper system.
  \end{abstract}

  \vspace*{1cm}
\end{center}

\newpage
%%

%% Table of Contents
\tableofcontents
\newpage
%%

\section{Documentation edition and implementation scope}
This is the 5.3.0 documentation edition, dated 2026-11-09. It carries forward the system and service contracts from the retained 5.2.2 documentation and identifies the current primary implementation sources. The former \path|documentation/v5.2.2| folder is retained unchanged.

Document version 5.3.0 is distinct from executable, dependency, registered system/service and wire-contract versions. Values still defined by the source (including 5.2.2 artefact versions, 1.0.0 service versions and \path|SysMonitor-AMQP-IDD-5.2.2| payload identifiers) are preserved. This documentation update does not change runtime behaviour or establish a new implementation release. Historical verification statements describe the tests recorded in the inherited baseline; this revision validates document compilation and references.

\section{Overview}
\label{sec:overview}

This document describes the \texttt{global\allowbreak{}Service\allowbreak{}Discovery} service interface of the Gatekeeper system, which provides discovery of authorized providers of a service in the neighbour clouds, through the inter-cloud relay.

It's implemented using protocol, encoding as stated in below table:
\begin{table}[ht!]
  \centering
  \begin{tabular}{|l|l|l|}
    \hline
    \rowcolor{gray!33} Profile type & Type & Version \\ \hline
    Transfer protocol & HTTP & 1.1 \cite{rfc9112} \\ \hline
    Data encryption & TLS & 1.3 \cite{rfc8446}, 1.2 \cite{rfc5246} \\ \hline
    Encoding & JSON & RFC 8259 \cite{rfc8259} \\ \hline
    Compression & N/A & - \\ \hline
    Semantics & N/A & - \\ \hline
    Ontology & N/A & - \\ \hline
  \end{tabular}
  \caption{Communication and semantics details used for the \texttt{global\allowbreak{}Service\allowbreak{}Discovery} service interface}
  \label{tab:comunication_semantics_profile}
\end{table}

This document provides the Interface Design Description IDD to the \textit{globalServiceDiscovery -- Service Description} document \cite{gk-globalServiceDiscovery-sd}.
For further details about how this service is meant to be used, please consult that document.

The rest of this document describes how to realize the \texttt{global\allowbreak{}Service\allowbreak{}Discovery} service
interface in detail. Both in terms of its operations (Section \ref{sec:operations}) and its information model (Section \ref{sec:model}).

\subsection{Interface registration}

The service is registered in the ServiceRegistry with the interface template \texttt{generic\_\allowbreak{}https}, policy \texttt{CERT\_\allowbreak{}AUTH}, access address \texttt{gatekeeper}, port 8449 and base path \texttt{/\allowbreak{}gatekeeper}. The \texttt{operations} property of the interface lists the operations below with their method and path relative to the base path (path parameters are placeholders such as \texttt{\{id\}} that the consumer fills in):

\noindent\begin{longtable}{| p{5.2cm} | p{1.6cm} | p{8.4cm} |} \hline
\rowcolor{gray!33} Operation & Method & Path \\ \hline \endhead
\texttt{query} & POST & \texttt{/\allowbreak{}init\_\allowbreak{}gsd} \\ \hline
\texttt{multi-\allowbreak{}query} & POST & \texttt{/\allowbreak{}init\_\allowbreak{}multi\_\allowbreak{}gsd} \\ \hline
\end{longtable}

\newpage

\section{Service Operations}
\label{sec:operations}

The interfaces of the \texttt{global\allowbreak{}Service\allowbreak{}Discovery} service, its operations, data models and implementation are provided below.
In particular, each subsection first names the HTTP method and path used to call the operation, after which it names the abstract operation from the \texttt{global\allowbreak{}Service\allowbreak{}Discovery} SD document, as well as input and output types.

Every request needs mutual TLS with a system or operator certificate of the cloud and \texttt{Authorization: Bearer IDENTITY-TOKEN//<token>}; the token is verified with Authentication and must belong to the system named by the certificate. Request and response bodies are \texttt{application/json} unless stated otherwise.
Errors are answered with \texttt{\{\"error\": \"<text>\"\}} and the status codes \texttt{400 Bad Request} (invalid input), \texttt{401 Unauthorized} (no valid identity), \texttt{403 Forbidden} (caller not permitted), \texttt{404 Not Found}, \texttt{500 Internal Server Error}, \texttt{502 Bad Gateway} (a system it depends on failed) and \texttt{503 Service Unavailable} (overloaded).

\osubsection{query}{POST}{/\allowbreak{}gatekeeper/\allowbreak{}init\_gsd}{\mref[Global\allowbreak{}Service\allowbreak{}Discovery\allowbreak{}Request]{GlobalServiceDiscoveryRequest}}{\mref[Global\allowbreak{}Service\allowbreak{}Discovery\allowbreak{}Result]{GlobalServiceDiscoveryResult}}

Asks the target clouds for authorized providers of the required service.

A successful call returns \texttt{200 OK} with the output in the body.

\begin{lstlisting}[label={lst:query},caption={An \oref{query} invocation.}]
POST /gatekeeper/init_gsd HTTP/1.1
Host: gatekeeper:8449
Authorization: Bearer IDENTITY-TOKEN//<token>
Content-Type: application/json

{
  "serviceRequirement": {
    "serviceDefinition": "temperatureReading",
    "operations": [
      "read"
    ]
  },
  "targetClouds": [
    "GoCloud|Company"
  ],
  "consumerSystem": "TemperatureDisplay"
}
\end{lstlisting}

\begin{lstlisting}[label={lst:query-resp},caption={An \oref{query} response.}]
{
  "results": [
    {
      "cloud": "GoCloud|Company",
      "providerName": "TemperatureSensor",
      "serviceDefinition": "temperatureReading",
      "version": "5.2.2",
      "metadata": {
        "unit": "celsius"
      },
      "operations": [
        "read"
      ]
    }
  ],
  "unsuccessfulRequests": 0,
  "queriedClouds": 1
}
\end{lstlisting}

\osubsection{multi-query}{POST}{/\allowbreak{}gatekeeper/\allowbreak{}init\_multi\_gsd}{\mref[Global\allowbreak{}Service\allowbreak{}Discovery\allowbreak{}Request]{GlobalServiceDiscoveryRequest}}{\mref[Global\allowbreak{}Service\allowbreak{}Discovery\allowbreak{}Result]{GlobalServiceDiscoveryResult}}

Same as \texttt{query} (v4 multi-GSD path).

A successful call returns \texttt{200 OK} with the output in the body.

\begin{lstlisting}[label={lst:multi-query},caption={An \oref{multi-query} invocation.}]
POST /gatekeeper/init_multi_gsd HTTP/1.1
Host: gatekeeper:8449
Authorization: Bearer IDENTITY-TOKEN//<token>
Content-Type: application/json

{
  "serviceRequirement": {
    "serviceDefinition": "temperatureReading",
    "operations": [
      "read"
    ]
  },
  "targetClouds": [
    "GoCloud|Company"
  ],
  "consumerSystem": "TemperatureDisplay"
}
\end{lstlisting}

\begin{lstlisting}[label={lst:multi-query-resp},caption={An \oref{multi-query} response.}]
{
  "results": [
    {
      "cloud": "GoCloud|Company",
      "providerName": "TemperatureSensor",
      "serviceDefinition": "temperatureReading",
      "version": "5.2.2",
      "metadata": {
        "unit": "celsius"
      },
      "operations": [
        "read"
      ]
    }
  ],
  "unsuccessfulRequests": 0,
  "queriedClouds": 1
}
\end{lstlisting}

\newpage

\section{Data Models}
\label{sec:model}

Here, all data objects that can be part of the service calls associated with this service are listed in alphabetic order.
Note that each subsection, which describes one type of object, begins with the \textit{struct} keyword, which is meant to denote a JSON \pref{Object} that must contain certain fields, or names, with values conforming to explicitly named types.
Fields that are not needed may be omitted or \texttt{null} in requests; empty fields are omitted in responses.
As a complement to the primary types defined in this section, there is also a list of secondary types in Section \ref{sec:model:primitives}, which are used to represent things like hashes, identifiers and texts.

\msubsection{struct}{GlobalServiceDiscoveryRequest}




\noindent\begin{tabularx}{\textwidth}{| >{\raggedright\arraybackslash}p{3.9cm} | >{\raggedright\arraybackslash\small}p{5.3cm} | X |} \hline
\rowcolor{gray!33} Object Field & Value Type & Description \\ \hline
"service\allowbreak{}Requirement" & \pref{Object} & Orchestration service requirement; \texttt{serviceDefinition} is required, \texttt{operations} are checked (alias \texttt{requestedService}). \\ \hline
"target\allowbreak{}Clouds" & \pref{Array}$<$\allowbreak{}\pref{String}$>$ & Clouds to ask (\texttt{\textless{}Cloud\textgreater{}} or \texttt{\textless{}Cloud\textgreater{}\textbar{}\allowbreak{}\textless{}Organization\textgreater{}}); default all configured peer clouds. \\ \hline
"consumer\allowbreak{}System" & \pref{String} & Consumer on whose behalf the orchestrator asks; default the requester. \\ \hline
\end{tabularx}



\msubsection{struct}{GlobalServiceDiscoveryResult}




\noindent\begin{tabularx}{\textwidth}{| >{\raggedright\arraybackslash}p{3.9cm} | >{\raggedright\arraybackslash\small}p{5.3cm} | X |} \hline
\rowcolor{gray!33} Object Field & Value Type & Description \\ \hline
"results" & \pref{Array}$<$\allowbreak{}\mref{InterCloudProvider}$>$ & Authorized providers of all asked clouds. \\ \hline
"unsuccessful\allowbreak{}Requests" & \pref{Number} & Number of clouds that did not answer correctly. \\ \hline
"queried\allowbreak{}Clouds" & \pref{Number} & Number of clouds asked. \\ \hline
\end{tabularx}



\msubsection{struct}{InterCloudProvider}




\noindent\begin{tabularx}{\textwidth}{| >{\raggedright\arraybackslash}p{3.9cm} | >{\raggedright\arraybackslash\small}p{5.3cm} | X |} \hline
\rowcolor{gray!33} Object Field & Value Type & Description \\ \hline
"cloud" & \pref{String} & Provider cloud, \texttt{\textless{}Cloud\textgreater{}\textbar{}\allowbreak{}\textless{}Organization\textgreater{}}. \\ \hline
"provider\allowbreak{}Name" & \pref{String} & Provider system. \\ \hline
"service\allowbreak{}Definition" & \pref{String} & Service definition. \\ \hline
"version" & \pref{Version} & Service version. \\ \hline
"metadata" & \pref{Object} & Instance metadata. \\ \hline
"operations" & \pref{Array}$<$\allowbreak{}\pref{String}$>$ & Offered operations. \\ \hline
\end{tabularx}



\subsection{Primitives}
\label{sec:model:primitives}

As all messages are encoded using the JSON format \cite{rfc8259}, the following primitive constructs, part of that standard, become available.
Note that the official standard is defined in terms of parsing rules, while this list only concerns syntactic information.
Furthermore, the \pref{Object} and \pref{Array} types are given optional generic type parameters, which are used in this document to signify when pair values or elements are expected to conform to certain types.

\noindent\begin{tabularx}{\textwidth}{| p{4.6cm} | X |} \hline
\rowcolor{gray!33} Type & Description \\ \hline
\pdef{Array} & A JSON Array \cite{rfc8259}; \pref{Array}$<$A$>$ has elements of type A. \\ \hline
\pdef{Number} & JSON Number, an integer. \\ \hline
\pdef{Object} & JSON Object; \pref{Object}$<$A$>$ has string keys and values of type A. \\ \hline
\pdef{String} & JSON String. \\ \hline
\pdef{Version} & JSON String, semantic version \texttt{major.minor.patch}. \\ \hline
\end{tabularx}

\newpage

\section{Go and Rust implementation variants}
\subsection{Go implementation}
Primary system folder: \path|Eclipse-Arrowhead-go/ah5-Gatekeeper-go|.\par
The system owns its Go module, launcher and build configuration. Use the per-system \path|Makefile| and \path|go.mod|; the containing Go tree is the build source, not a historical Arrowhead copy. Ordinary native roles use independent identities and JSON snapshot state rather than the Java database. Do not assume Java SQL migration, clustering or identical transaction/retry semantics. The registry has its own configuration and optional native AMQP monitor; other roles must be checked against their own registrations.
\paragraph{Listening port.} Built-in default 127.0.0.1:8449 (Arrowhead default). The shipped config.json, compose.yaml and Dockerfile use the same port on all interfaces ("listen": "0.0.0.0:8449"; 8540--8555 layout before 2026-11-09). Override with "listen" in the configuration file, SERVER\_PORT=N or -server.port N; the port overrides keep the configured host.
\paragraph{Native behavior and outstanding work.} Native Go legacy-mTLS cloud discovery/negotiation engine adapted to this runtime. AH-GT-1 encryption, complete current Java edge-case parity, real peer/broker/ACL qualification. These are the boundaries recorded by the native compatibility document; later implementation changes still require direct qualification.
Native design/route evidence: \path|Eclipse-Arrowhead-go/ah5-Gatekeeper-go/documentation/Gatekeeper|.\par
For this Java interface, select only a native transport and operation explicitly implemented and advertised by that port. A matching method/path does not establish DTO, error, authorization, retry or persistence equivalence. Java MQTT/AMQP bindings are not automatically provided by a Go or Rust role. Differential contract and security checks remain necessary before substitution.
\subsection{Rust implementation}
Primary system folder: \path|Eclipse-Arrowhead-rust/ah5-Gatekeeper-rust|.\par
The per-system entry point is declared as a binary of \path|arrowhead-support-rust| in \path|ah5-support-runtime-rust/Cargo.toml|. Build from the primary Rust workspace with the applicable package/binary selection. Documentation links can resolve into the shared core/support catalogue. Independent process identity and state remain required.\par
The ordinary native profile uses Tokio/Axum and rustls with peer-identity checks and independent locked JSON state. Network effects are not distributed database transactions. The precise trust, CRL reload and persistence guarantees are those of the native profile, not the Java implementation.\par
\paragraph{Listening port.} Default 127.0.0.1:8449 (Arrowhead default, role gatekeeper). Override with "listen": "host:port" in the configuration, SERVER\_PORT=N or --server.port N; the port overrides keep the configured host.
\paragraph{Native behavior and outstanding work.} Neighbour/relay administration, actual MQTT global discovery and negotiation with explicit cloud/organisation bindings. Exhaustive DTO/filter/failure parity and independent review. These are the boundaries recorded by the native compatibility document; later implementation changes still require direct qualification.
Native design/route evidence: \path|Eclipse-Arrowhead-rust/ah5-support-runtime-rust/documentation/Gatekeeper|.\par
For this Java interface, select only a native transport and operation explicitly implemented and advertised by that port. A matching method/path does not establish DTO, error, authorization, retry or persistence equivalence. Java MQTT/AMQP bindings are not automatically provided by a Go or Rust role. Differential contract and security checks remain necessary before substitution.

\bibliographystyle{IEEEtran}
\bibliography{bibliography}

\newpage

\section{Revision History}
\subsection{Amendments}

\noindent\begin{tabularx}{\textwidth}{| p{1cm} | p{3cm} | p{2cm} | X | p{4cm} |} \hline
\rowcolor{gray!33} No. & Date & Version & Subject of Amendments & Author \\ \hline

1 & 2026-10-22 & 5.2.2 & Document created from the Arrowhead LaTeX templates for the 5.2.2 Java Spring implementation, including the local extensions of this workspace. & Jerker Delsing \\ \hline
2 & 2026-10-30 & 5.2.2 & Checked against the primary Java source tree (Eclipse-Arrowhead-java): services, operations, paths and data models unchanged. & Claude Opus 5.5 (draft) \\ \hline
3 & 2026-11-09 & 5.3.0 & New documentation edition; primary-source references and Go/Rust implementation scope; 5.2.2 retained. & Documentation maintenance \\ \hline
\end{tabularx}

\subsection{Quality Assurance}

\noindent\begin{tabularx}{\textwidth}{| p{1cm} | p{3cm} | p{2cm} | X |} \hline
\rowcolor{gray!33} No. & Date & Version & Approved by \\ \hline

1 & & \arrowversion  &  \\ \hline

\end{tabularx}

\end{document}
